\section{Introduction}\label{sec:intro}

Serving large language models is now one of the largest uses of data-centre GPUs~\cite{kessler2025inference}.
Operators who want to report the carbon footprint of a request, or price it fairly, need to know how much energy that request used.
On a dedicated GPU the answer is the integral of the board's power over the request's runtime.
Production deployments rarely look like that.
Requests are batched, models share devices, and the power sensor reports one number for everything that runs on the board~\cite{varga2024sensors}.

Existing tools attribute energy by time share~\cite{quintero2023carbon}, or they assume that every request in a batch costs the same.
Both assumptions break down under load.
In our measurements a short prompt that joins a full batch costs a third of what the same prompt costs on an idle GPU, while a long generation in a busy batch costs more than the time-share estimate suggests (Section~\ref{sec:eval}).

This paper makes three contributions.
\begin{itemize}
  \item \sysname{}, a method that splits GPU energy across concurrent inference requests using kernel traces and the device's power readings (Section~\ref{sec:method}).
  \item An evaluation on four open-weight models and two GPU generations against isolated reference runs (Section~\ref{sec:eval}).
  \item A public data set of 1.2 million attributed requests.
\end{itemize}
